\documentclass[11pt,a4paper]{article}
\usepackage[margin=0.85in,top=1in,bottom=1in]{geometry}
\usepackage{amsmath,amssymb,amsfonts}
\usepackage{graphicx}
\usepackage{float}
\usepackage{enumitem}
\usepackage{microtype}
\usepackage{caption}
\usepackage[hidelinks]{hyperref}
\usepackage[most,breakable]{tcolorbox}
\tcbuselibrary{breakable,skins}
\usepackage[utf8]{inputenc}
\usepackage[T1]{fontenc}
\usepackage{lmodern}
\captionsetup{font=small,labelfont=bf,skip=4pt}
\setlist{leftmargin=1.5em,topsep=2pt}

%% Breakable card — prevents content cutoff across pages
\newtcolorbox{solcard}[3]{
  breakable, enhanced,
  colback=white, colframe=gray!50, arc=2pt,
  title={\textbf{#1}\hfill\textsf{\small #3\,pts}},
  fonttitle=\normalsize\sffamily\bfseries,
  before upper={\small\textit{#2}\par\smallskip\hrule height 0.4pt\smallskip},
  top=4pt, bottom=6pt, left=7pt, right=7pt,
  parbox=false,
}

\title{\textbf{Champagne foam --- Solution}}
\author{\normalsize X22 (2022) — English translation}
\date{}
\begin{document}\maketitle\vspace{-0.5em}
\begin{solcard}{A1}{Let the total number of moles of carbon dioxide gas in a closed bottle be $\nu_0 = 0.05$ mol, the volume of liquid $V_L =250$ ml, and the volume of gas $V_G = 5$ ml. Find the pressure $P$ of carbon dioxide in the bottle. Please also provide a numerical value. \textbf{In further paragraphs, assume that the liquid is in an open container.}}{0.80}

Let $\nu_G$ be the amount of carbon dioxide substance in gaseous form, $\nu_L$ – in solution. The amount of substance in a gas can be expressed in terms of pressure using the ideal gas equation of state. The concentration of carbon dioxide in a solution is expressed in terms of pressure using Henry's law, then we get the equation $$ \nu_0 = \nu _G + \nu _L = \frac{PV_G}{RT} + cV_L = P\left( \frac{V_G}{RT} + k_H V_L\right), $$ from which pressure $$ P = \frac{\nu_0 R T}{V_G + k_H V_L RT}. $$

\par\smallskip\noindent\textbf{Answer:} $$ P = \frac{\nu_0 R T}{V_G + k_H V_L RT} \approx 5.4 \cdot 10^5~\text{Pa} $$
\end{solcard}

\begin{solcard}{A2}{A bubble with carbon dioxide can form only if the pressure of carbon dioxide in it can balance the pressure created by surface tension forces. Find the minimum bubble radius $R^*$ that can form. Express the answer in terms of the molar concentration of gas in the liquid $c$, atmospheric pressure $P_0$, as well as physical constants characterizing the liquid. What is the concentration $c$ if the critical radius is $R^* = 1.5 \cdot 10^{-6}$ m?}{0.70}

The pressure of carbon dioxide in the bubble must balance the external pressure (which we consider equal to atmospheric pressure $P_0$) and the addition due to surface tension $$ \frac{2 \gamma}{R^*} + P_0 = P_{CO_2} = \frac{c}{k_H}. $$ U bubble one surface, therefore surface tension creates pressure $2 \gamma/R$. Solve equation, obtain $$ \frac{2 \gamma}{R^*} = \frac{c - k_H P_0}{k_H},\quad R^* = \frac{2 \gamma k_H}{c - k_H P_0}. $$ Concentration one can calculate by formula $$ c = k_H\left( P_0+\frac{2 \gamma}{R^*}\right) \approx  61 \frac{\text{mol}}{\text{m}^3}. $$

\par\smallskip\noindent\textbf{Answer:} $$ R^* = \frac{2 \gamma k_H}{c - k_H P_0} $$ $$ c = k_H\left( P_0+\frac{2 \gamma}{R^*}\right) \approx  61 \frac{\text{mol}}{\text{m}^3} $$
\end{solcard}

\begin{solcard}{A3}{Using the above considerations, find the diffusion coefficient $B$. Express your answer in terms of $k_B$, $T$, $\beta$. Determine the mobility of molecules $\beta$ in terms of viscosity $\eta$ and the size of the molecules. Find the numerical value of the diffusion coefficient.  \textbf{If you fail to solve this item, you can use the value} $B = 5 \cdot 10^{-9} \text{m}^2/\text{s}$ to obtain numerical answers in the future.}{1.00}

The average speed of a particle under the action of force $F_z=F$ is $v= \beta F$, so the force creates a flow $j_z^F = n v = n \beta F$. The concentration depends on the coordinate as:

\par\smallskip\noindent\textbf{Answer:} $$ B = \frac{k_B T}{6 \pi \eta r} \approx 1.4 \cdot 10^{-9} \mathrm{m}^2/c $$
\end{solcard}

\begin{solcard}{B1}{Find how the diameter $D$ of the bubble depends on time. Consider the initial diameter to be approximately equal to zero, express the answer through $f_1$, $f_2$, $c$, $c_i$, $R$, $T$, $P_0$. The pressure in the bubble and temperature can be considered constant; the contribution of surface tension and pressure of the liquid column to the gas pressure in the bubble can be neglected.}{1.00}

Let us express the volume of the bubble through the amount of substance in it $$ V = \frac{RT}{P} \nu. $$ From the formula from the condition obtain velocity change volume $$ \frac{dV}{dt} = \frac{RT}{P} \frac{d\nu}{dt} = \frac{RT}{P}  k S (c-c_i). $$ Substitute expression through diameter: $$ 3 D^2 f_2\frac{dD}{dt}  = \frac{RT}{P}  \frac{B}{\kappa D} f_1 D^2 (c-c_i) $$ Hence we get equation $$ 2D\frac{dD}{dt} =  \frac{2f_1}{3\kappa f_2} \frac{RT}{P}B(c-c_i), $$ on the left side of which is the derivative of the square of the diameter with respect to time.

\par\smallskip\noindent\textbf{Answer:} $$ D  =  \sqrt{\frac{2f_1}{3\kappa f_2} \frac{RT}{P}B(c-c_i) t} $$
\end{solcard}

\begin{solcard}{B2}{Find expressions for $f_1(\theta)$ and $f_2(\theta)$. Construct a (qualitative) graph of the ratio $f_1/f_2$ to $\theta \in (0,\,\pi/2)$.}{0.70}

The surface area of the sphere is equal to the product of the square of the radius $D^2/4$ and the solid angle at which it is visible from the center, therefore $$ f_1 = \frac{\pi }{2} \left(1 + \cos \theta  \right) $$ Volume part sphere one can find as sum volume of the spherical sector, visible under the more same solid angle,  and cone. Obtain $$ f_2 = \frac{\pi}{24}\left(2 + 3\cos \theta -  \cos^3 \theta \right) $$

\par\smallskip\noindent\textbf{Answer:} $$ f_1 = \frac{\pi }{2} \left(1 + \cos \theta  \right) $$ $$ f_2 = \frac{\pi}{24}\left(2 + 3\cos \theta -  \cos^3 \theta \right) $$
\begin{figure}[H]
  \centering
  \includegraphics[width=0.72\textwidth]{T5_S_EN_assets/fig1.jpg}
\end{figure}
\end{solcard}

\begin{solcard}{B3}{Find the bubble radius $R$ at which it comes off the surface. Consider that it is attached to the surface along a circle of radius $a_d \ll R$, the angle between the bubble and the surface is $\theta$,}{0.30}

The surface tension force must balance the Archimedes force $$ 2 \pi a_d \gamma \sin \theta = \frac{4 \pi}{3} \rho g R^3 $$ acting on the bubble

\par\smallskip\noindent\textbf{Answer:} $$ R = \left(\frac{3 a_d \gamma \sin \theta}{2\rho g} \right)^{1/3} $$
\end{solcard}

\begin{solcard}{B4}{After the bubble breaks away from the surface, a region of liquid with a reduced concentration of dissolved carbon dioxide remains near it. Because of this, some time $t_n$ must pass before the next bubble begins to form in the same place. Consider that the concentration of carbon dioxide is reduced in the area of size $D_m$, and the depth of the defect that needs to be filled with carbon dioxide is $h \ll D_m$. The dependence of concentration on distance to the surface is linear. Find the time $t_n$ during which the bubble that appears at the bottom of the defect grows to the surface of the glass. The concentration near the defect is equal to the minimum concentration $c_n$ at which the existence of a bubble is possible; the dependence of the concentration on distance can be considered linear.}{0.80}

Due to the difference in concentrations, a flow of carbon dioxide $$ j = B \frac{c - c_n}{D_m}. $$ Note, that one and that same coefficient diffusion connects as flux/flow molecule with difference concentration, thus and molar flux/flow with difference molar concentration, since ordinary and molar concentration proportional to each other. Let area defectand $S$. Volume bubble relate with amount substance in it $$ V = \nu \frac{RT}{P}. $$ Velocity change amount substance in bubble $$ \frac{d\nu}{dt} = jS= \frac{c-c_n}{D_m}BS = \frac{P}{RT}\frac{dV}{dt} = \frac{PS}{RT} \frac{dh}{dt} $$ Hence it follows, that defect is filled with constant velocity $$ \frac{dh}{dt} = \frac{RT}{P} B\frac{c - c_n}{D_m}, $$ and time formation bubble $$ t_n =  \frac{P}{RT} \frac{h D_m}{B(c-c_n)}. $$ appears

\par\smallskip\noindent\textbf{Answer:} $$ t_n =  \frac{P}{RT} \frac{h D_m}{B(c-c_n)} $$
\end{solcard}

\begin{solcard}{B5}{Let $t_g$ be the time during which the bubble grows from the edge of the defect to its maximum diameter $D_m$. When the concentration of carbon dioxide in the liquid changes, this time, as well as the time of bubble formation $t_n$, can change. In what coordinates will the dependence of $t_g$ on $t_n$ be linear?}{0.70}

The time for the bubble to grow to its maximum diameter can be found from the results of point B1: $$ t_g = \frac{D_m^2}{K_1(c-c_i)}, $$ where $K_1$ – some constant, not depend from concentration. $D_m$ be determined in main equilibrium force surface tension and force Archimedes, therefore also not depend from concentration. Then dependence inverse time from concentration carbon gas in solution linear: $$ \frac{1}{t_g} \sim c-c_i, \quad \frac{1}{t_n} \sim c - c_n. $$ Therefore, the relationship between the inverse times themselves should also be linear.

\par\smallskip\noindent\textbf{Answer:} Dependence in coordinates $1/t_g$ on $1/t_n$.
\begin{figure}[H]
  \centering
  \includegraphics[width=0.72\textwidth]{T5_S_EN_assets/fig2.jpg}
\end{figure}
\end{solcard}

\begin{solcard}{C1}{The rate at which a bubble floats up is determined by the force of viscous friction acting on it, which can be calculated using the Stokes formula $F = 6 \pi \eta aU$, where $a$ is the radius of the bubble, $U$ is the velocity of the bubble, $\eta$ is the viscosity of the liquid. Find the rising speed of a bubble of radius $a$.}{0.40}

We believe that the speed of the bubble is established quite quickly. Then it is determined by the equilibrium between the force of viscous friction and the Archimedes force $$ \rho g \frac{4 \pi}{3} a^3 = 6 \pi \eta a U. $$ Hence velocity $$ U = \frac{2 \rho g}{9 \eta}a^2. $$

\par\smallskip\noindent\textbf{Answer:} $$ U = \frac{2 \rho g}{9 \eta}a^2 $$
\end{solcard}

\begin{solcard}{C2}{Find the rate of change of the radius of the bubble $da/dt$. Express your answer in terms of $a$, concentration difference $\Delta c$, temperature $T$, atmospheric pressure $P_0$, $\Delta c$ and physical constants.}{0.60}

Let's find the value of the coefficient $K$ for the bubble: $$ K = 0.63 B^{2/3} \frac{U^{1/3}}{a^{2/3}} = 0.63 B^{2/3} \left(\frac{2 \rho g}{9 \eta}\right)^{1/3} =  0.38 B^{2/3} \left(\frac{ \rho g}{ \eta}\right)^{1/3}. $$ Turn out, it not depend from radius. Similarly part B relate express velocity change volume: $$ \frac{dV}{dt} = \frac{RT}{P} \frac{d\nu}{dt} =  \frac{RT}{P} KS \Delta c. $$ Change volume relate with change radius relation $$ \frac{dV}{dt} =  S \frac{da}{dt}  = \frac{RT}{P} KS \Delta c, $$ from which $$ \frac{da}{dt} =  0.38 B^{2/3} \left(\frac{ \rho g}{ \eta}\right)^{1/3} \frac{RT}{P} \Delta c = v. $$ This speed does not depend on the radius of the bubble.

\par\smallskip\noindent\textbf{Answer:} $$ \frac{da}{dt} =  0.38 B^{2/3} \left(\frac{ \rho g}{ \eta}\right)^{1/3} \frac{RT}{P} \Delta c $$
\end{solcard}

\begin{solcard}{C3}{The bubble floats up from the bottom of the vessel, initial radius $a_0$. Find its radius $a$ and volume $V_b$ near the surface.}{0.80}

Let us write the equation for the growth of the bubble radius $$ \frac{da}{dt} = v, $$ and also for motion bubble ($h$ – height lift) $$ \frac{d h}{dt} = U = \frac{2 \rho g}{9 \eta}a^2. $$ Hence $$ \frac{da}{dh} = \frac{1}{a^2} \frac{9 \eta v}{2 \rho g}. $$ Integrate, obtain $$ a^3 = \frac{27 \eta v}{2 \rho g} h + a^3_0. $$

\par\smallskip\noindent\textbf{Answer:} $$ a = \left(\frac{27 \eta v}{2 \rho g} h + a^3_0\right)^{1/3} $$ $$ V_b = \frac{4 \pi}{3} \left(\frac{27 \eta v}{2 \rho g} h + a^3_0\right) $$ where $v =0.38 B^{2/3} \left(\frac{ \rho g}{ \eta}\right)^{1/3} \frac{RT}{P} \Delta c$
\end{solcard}

\begin{solcard}{C4}{Bubbles will form until the concentration of carbon dioxide in the liquid drops from the initial value $c$ to some critical value $c_n$, at which bubbles will no longer form. Find the total number of bubbles $N$ that can appear in the entire time. Express the answer in terms of the volume of the liquid $V_L$, its depth $h$, the initial and final concentrations of carbon dioxide and the constants characterizing the liquid and gas. Find also the numerical value. Use data $c=100~\text{mol}/\text{m}^3$, $c_n = 50~\text{mol}/\text{m}^3 $, $V_L = 0.5$ liters, vessel diameter $\xi = 6$ cm. Consider that the concentration near the surface of the bubble is such that the carbon dioxide gas in the bubbles is in equilibrium with that dissolved in the liquid. All bubbles float up from the bottom of the vessel, and their initial radius is zero. The change in concentration during the time from the formation of the bubble until it reaches the surface can be neglected.}{1.50}

With each bubble, an amount of carbon dioxide substance comes out of the liquid $$ \nu_b = 10.7 \left( \frac{B \eta}{\rho g}\right)^{2/3} h \Delta c. $$ It one can obtain, multiply volume bubble on $P/RT$. Difference concentration near bubble and in volume liquid $$ \Delta c = c - k_hP_0 $$ Change concentration carbon gas in liquid for output one  bubble $$ dc =- \frac{\nu _b}{V_L}. $$ Number released bubble relate with change concentration relation $$ dN = -\frac{dc V_L}{\nu_b} =- 0.093\left( \frac{\rho g }{B \eta} \right)^{2/3} \frac{V_L}{h}\frac{dc}{c - k_H P_0} $$ Integrate, find number bubble $$ N = 0.093\left( \frac{\rho g }{B \eta} \right)^{2/3} \frac{V_L}{h} \ln \frac{c - k_H P_0}{c_n - k_H P_0} \approx 1.1 \cdot 10^7 $$ Numerical value lead to for correct coefficient diffusion. For approximate value $5 \cdot 10^{-9} ~m^2/s$ obtain $N \approx 4.6 \cdot 10^6$.

\par\smallskip\noindent\textbf{Answer:} $$ N = 0.093\left( \frac{\rho g }{B \eta} \right)^{2/3} \frac{\pi \xi^2}{4} \ln \frac{c - k_H P_0}{c_n - k_H P_0} \approx 1.1 \cdot 10^7 $$
\end{solcard}

\begin{solcard}{C5}{Let us assume that all the bubbles formed in the process are not destroyed, and their sizes are equal to the sizes of the bubbles reaching the surface of the liquid at a finite concentration of carbon dioxide. What height of foam $H$ is formed on the surface of the liquid?}{0.70}

The height of the foam is related to the volume of bubbles by the relation $H = V/S$, where the volume of foam is $V=V_b N$. The bubble volume corresponding to the final concentration value $$ V_b = 10.7 \left( \frac{B \eta}{\rho g}\right)^{2/3} h (c_n - k_H P_0) \frac{RT}{P}, $$ volume foam $$ V = \frac{\pi \xi^2}{4} h(c_n - k_H P_0) \frac{RT}{P} \ln \frac{c - k_H P_0}{c_n - k_H P_0} , $$ and means it height $$ H = \frac{RT}{P} (c_n - k_H P_0) h \ln \frac{c - k_H P_0}{c_n - k_H P_0} \approx 0.53 h. $$ Here is used height column liquid $h = V/(\pi \xi^2/4) \approx 17.7 ~cm $.

\par\smallskip\noindent\textbf{Answer:} $$ H = \frac{RT}{P} (c_n - k_H P_0) h \ln \frac{c - k_H P_0}{c_n - k_H P_0} \approx 0.53 h \approx 9 ~cm $$
\end{solcard}

\end{document}